\section{Введение}
\label{sec:Chapter0} \index{Chapter0}

В последние годы практически индустриальным стандартом разработки и эксплуатации программного обеспечения стало использование \textit{контейнерных технологий}. Контейнеризация позволяет упаковать приложение вместе со всеми его зависимостями в изолированный артефакт, исполнение которого не зависит от особенностей окружения хоста и не требует от него ничего, кроме поддержки исполнения контейнеров. Стандартное представление таких артефактов закреплено в спецификациях \textit{Open Container Initiative (OCI)}~\cite{ociimagespec, ocidistspec, ociruntimespec}, благодаря чему различные инструменты сборки и исполнения контейнеров взаимозаменяемы. Контейнерные образы стали универсальной единицей доставки программного обеспечения, что особенно ярко проявляется в облачной инфраструктуре~--- например, при использовании Kubernetes, serverless-вычислений или CI/CD-конвейеров.

Однако с распространением контейнеризации возникает и новая поверхность атаки: контейнер обычно содержит не только код приложения, но и большое количество стороннего программного обеспечения~--- пакеты Linux-дистрибутива, разделяемые библиотеки, языковые зависимости. Все они могут содержать известные уязвимости (\textit{CVE})~\cite{cve}, и их наличие напрямую отражается на безопасности конечного приложения. Учитывая темпы пополнения соответствующих баз данных~\cite{nvd}, регулярная и автоматизированная проверка образов на наличие уязвимостей становится обязательной частью цикла поставки программного обеспечения. Эта задача традиционно возлагается на инструменты, называемые \textit{сканерами уязвимостей}~--- такими как Trivy~\cite{trivy}, Clair~\cite{clair} и Grype~\cite{grype}.

В организациях, где контейнерные образы создаются и хранятся централизованно, эта задача обычно решается путём интеграции сканера с \textit{container registry}~--- сервисом, отвечающим за хранение и распространение образов. Особенно характерно это для \textit{managed container registry} облачных провайдеров, а также для собственных корпоративных registry-инсталляций, например, на базе CNCF-проекта Harbor~\cite{harbor}. В таких сценариях сканирование производится со стороны сервера: каждый загружаемый образ автоматически проверяется на наличие уязвимостей, а результаты сохраняются и используются для последующих политик доступа и развертывания.

Несмотря на функциональную зрелость существующих сканеров, при их использовании в подобной серверной модели обнаруживается общий недостаток: для анализа образа сканер скачивает все его слои \textit{полностью}, тогда как для проверки на уязвимости в большинстве случаев требуется лишь небольшая часть содержимого~--- метаданные пакетных менеджеров операционной системы (списки установленных пакетов apt/apk/dnf) и манифесты языковых зависимостей (\texttt{requirements.txt}, \texttt{package.json}, \texttt{go.sum} и т.д.). Все остальные файлы~--- бинарные файлы программ, разделяемые библиотеки, файлы данных, кэши пакетных менеджеров~--- являются для сканера, как правило, \textit{избыточными}. Объемы таких <<лишних>> данных могут составлять подавляющую часть от полного размера образа, что напрямую сказывается на скорости сканирования и нагрузке на инфраструктуру container registry, особенно с учетом большого количества новых образов, регулярно загружаемых в registry в рамках CI/CD-процессов.

Уменьшение объема загружаемых из registry данных представляется хорошим вектором повышения эффективности сканирования. Одним из подходов к решению этой задачи является использование специальных форматов представления слоев образа, поддерживающих \textit{ленивую} (\textit{lazy}) загрузку~--- такой подход активно развивается в области ускорения запуска контейнеров и реализован, в частности, в формате \textit{eStargz}~\cite{estargzspec, stargzsnapshotter} и проектах семейства Nydus~\cite{nydus} и DADI~\cite{li2020dadi}. Особенно привлекательным с точки зрения адаптации в существующих сканерах представляется eStargz: он полностью совместим со стандартными форматами слоев OCI и не требует никаких изменений на стороне registry. Несмотря на наличие интеграций с такими форматами в составе container runtime, в сканерах уязвимостей эти возможности на момент написания работы не используются.

В данной работе исследуется возможность использования формата eStargz для повышения эффективности сканирования OCI-образов. На основе одного из наиболее популярных open-source сканеров Trivy~\cite{trivy} разрабатывается прототип, демонстрирующий применимость подхода, а также инструментарий для воспроизводимой экспериментальной оценки. Результаты подтверждают, что применение eStargz позволяет добиться существенного снижения объема загружаемых из registry данных и кратного ускорения сканирования по сравнению с обычным Trivy, при этом сохраняя корректность результатов и совместимость с существующей инфраструктурой.

Целями данной работы являются:
\begin{itemize}
\item изучение предметной области, в частности, способов представления и хранения OCI-образов, особенностей их использования в container registry, а также ключевых принципов работы сканеров уязвимостей;
\item обзор существующих сканеров уязвимостей с открытым исходным кодом и выявление общего недостатка, ограничивающего их эффективность при работе с registry;
\item предложение подхода к минимизации объема загружаемых из registry данных за счет использования формата eStargz, поддерживающего пофайловую загрузку содержимого слоев;
\item разработка прототипа сканера на базе Trivy, реализующего предложенный подход, а также инструментария для воспроизводимой экспериментальной оценки;
\item проведение экспериментальной оценки разработанного прототипа на тестовом наборе образов, моделирующем реальный профиль нагрузки на container registry.
\end{itemize}

\hyperref[sec:Chapter1]{\textit{Раздел 2}} посвящен обзору предметной области. Рассматривается устройство OCI-образов, принципы их хранения и распространения через container registry, а также описывается типичный сценарий интеграции сканеров уязвимостей в инфраструктуру.

В \hyperref[sec:Chapter2]{\textit{разделе 3}} разбираются ключевые принципы работы современных сканеров уязвимостей: модель сканирования, роль метаданных и пакетных манифестов, форматы представления результатов.

\hyperref[sec:Chapter3]{\textit{Раздел 4}} содержит обзор существующих сканеров с открытым исходным кодом (Trivy, Clair, Grype), анализ их работы с container registry, а также формулировку основной проблемы, решаемой в данной работе.

В \hyperref[sec:Chapter4]{\textit{разделе 5}} описываются предлагаемый подход к решению проблемы~--- ленивая загрузка слоёв через формат eStargz~--- и непосредственно реализация прототипа на базе сканера Trivy.

\hyperref[sec:Chapter5]{\textit{Раздел 6}} посвящен экспериментальной оценке: описывается методология бенчмарка, разработанный инструментарий и анализ полученных результатов в сравнении с существующими сканерами.

В \hyperref[sec:Chapter6]{\textit{разделе 7}} подводятся итоги работы, обозначаются ограничения предложенного подхода и обсуждаются возможные направления его развития.

\newpage
